\lysh{12857}{2}

\begin{alist}
\item 
\begin{rlist}
\item Για $\lambda = -2$ η εξίσωση γίνεται $(-2 - 1)x = 2 \cdot (-2) - 2 \Leftrightarrow -3x = -6 \Leftrightarrow x = 2$.
Άρα η λύση της εξίσωσης είναι $x = 2$.

\item Για $x = 1$ έχουμε:
\[(\lambda - 1) \cdot 1 = 2\lambda - 2 \Leftrightarrow \lambda - 1 = 2\lambda - 2 \Leftrightarrow \lambda = 1.\]
\end{rlist}

\item Η εξίσωση έχει άπειρες λύσεις αν και μόνο αν είναι της μορφής $0x = 0$, δηλαδή αν και μόνο αν 
\[\begin{cases} \lambda - 1 = 0 \\ 2\lambda - 2 = 0 \end{cases} \Leftrightarrow \lambda = 1.\]
\end{alist}
